\newpage\thispagestyle{hl}
\hlbanner{Part I $\cdot$ Architectural foundations}{2}{Mesh topologies and structural domain partitioning}
\hlsection{The question}
\begin{hlquote}A hand-written unit square passes the first check anyone would write, because its element areas add up to exactly 1. One of its elements is listed clockwise and gives a negative stiffness, and no error is raised.\end{hlquote}
\hlsection{Key concepts}
\begin{description}\setlength\itemsep{5pt}\setlength\parsep{0pt}\renewcommand{\descriptionlabel}[1]{\hspace{\labelsep}\sffamily\bfseries\color{hlnavy}#1}
\item[Connectivity] an integer array listing the nodes of each element; node numbers index the coordinate array
\item[Jacobian determinant] must be positive at every Gauss point; a negative sign marks an inverted element
\item[Graded mesh] element size set by a growth ratio between neighbours, not by a guessed exponent
\item[Conforming mesh] neighbouring elements share whole edges, with no hanging nodes
\item[Node set] a named selection of nodes, so supports and loads need no node numbers
\item[Mesh dispatcher] matches an element to the meshes it can run on
\end{description}
\hlsection{After this chapter you can}
\begin{itemize}\setlength\itemsep{3pt}
\item generate and grade a structured mesh, and join several element topologies in one mesh
\item check a mesh for positive Jacobians and for unshared interior edges
\item exchange a mesh through a file format and draw it, with a result, in two and three dimensions
\end{itemize}
\hlsection{In this chapter}
{\begin{multicols}{2}\raggedright
\textbf{2.1}\ What a mesh is\\
\textbf{2.2}\ Element topologies and node ordering\\
\textbf{2.3}\ Mesh quality\\
\textbf{2.4}\ Generators\\
\textbf{2.5}\ Grading\\
\textbf{2.6}\ Several element types in one mesh\\
\textbf{2.7}\ Building meshes from parts\\
\textbf{2.8}\ Boundary and node sets\\
\textbf{2.9}\ Exchange formats\\
\textbf{2.10}\ Visualising meshes and results\\
\textbf{2.11}\ The mesh dispatcher\\
\textbf{2.12}\ Summary
\end{multicols}}
\hlsection{Explained on the page, in short boxes}
\noindent Structured and unstructured meshes; Aspect ratio, angle and skewness; Connectivity order; Graded mesh and growth ratio; Conforming meshes and hanging nodes; Mesh file formats.
\newpage\thispagestyle{hl}
\noindent{\sffamily\small\color{hlgrey}Chapter 2 \quad Mesh topologies and structural domain partitioning}\par
\hlsection{Signature figure}
\hlfig{\hlimg{ch02_gallery2d}}{Structured meshes from the generators of \fea{}: rectangles, plates with holes, notches and fillets, with node and element counts.}
\hlsection{Headline result}
\begin{hlheadline}The clockwise element leaves the total area at 1.000, so only the Jacobian check reports it. An off-centre hole generator was found to give non-conforming meshes at four of six hole positions in the tagged release; it is conforming after the correction.\end{hlheadline}
\hlsection{From the code}
\begin{lstlisting}[style=hlpy]
good = rectangle_mesh(1.0, 1.0, 3, 3)          # a unit square, 9 elements
elems = good.elements.copy()
elems[4] = elems[4][::-1]                      # one element listed clockwise
bad = Mesh(good.nodes, elems, dim=2)
bad.check_quality(Quad4PlaneStress())
# -> *** 1 element(s) with non-positive detJ -- INVERTED/DEGENERATE ***
\end{lstlisting}
\hlsection{A stated limit}
\begin{hllimit}Structured meshes only: there is no unstructured generator or CAD import, and curved boundaries are polygons.\end{hllimit}
